% Generated draft: requires mathematical and source review.
\documentclass[11pt]{article}
\usepackage[margin=25mm]{geometry}
\usepackage{fontspec}
\setmainfont{Latin Modern Roman}
\usepackage{amsmath,amssymb,mathtools,unicode-math}
\setmathfont{Latin Modern Math}
\usepackage{enumitem,needspace,xurl,hyperref}
\setlength{\parindent}{0pt}
\setlength{\parskip}{5pt}
\setlength{\emergencystretch}{3em}
\newcommand{\N}{\mathbb{N}}
\newcommand{\Z}{\mathbb{Z}}
\newcommand{\Q}{\mathbb{Q}}
\newcommand{\R}{\mathbb{R}}
\newcommand{\C}{\mathbb{C}}
\newcommand{\F}{\mathbb{F}}
\newcommand{\A}{\mathbb{A}}
\newcommand{\PP}{\mathbb{P}}
\newcommand{\E}{\mathbb{E}}
\newcommand{\eps}{\varepsilon}
\newcommand{\dd}{\,\mathrm{d}}
\begin{document}
\section*{UnsolvedMath ID: 20002051}
\textbf{Source problem code:} catalogue entry
\subsection{Definitions and mathematical statement}
The reference to equation (40) is recovered from companion record 20002050. On a closed even-dimensional Riemannian manifold, let \(A_g\) be a natural elliptic differential operator with positive-definite leading symbol, assumed to be a positive integral power of a conformally covariant operator. The source examples include GJMS operators and the square of the Dirac operator. Fix the spectral and zero-mode conventions needed for its zeta-regularized determinant; write \(\det'\) when zero eigenvalues are omitted. For \(\widehat g=e^{2\omega}g\), the reported formula in dimensions 2, 4 and 6 has the form
\[
-\log\frac{\det' A_{\widehat g}}{\det' A_g}
=\alpha\left(\tfrac12\int_M\omega P_{n,g}\omega\,\dd V_g
+\int_M\omega Q_g\,\dd V_g\right)
+C_F(\widehat g)-C_F(g)+H(\widehat g,g),
\]
where \(C_F(g)=\int_MF_g\,\dd V_g\) is a local counterterm, \(P_n\) is the critical GJMS operator and \(H\) is a kernel cocycle, zero when the kernel is zero. A counterterm sign can be included in \(F\); hats and fractions are damaged in the extraction. Problem 29 asks whether an appropriate version holds in higher even dimensions, hence beyond the displayed 2, 4 and 6 cases.

The supplied aid reports a structural conditional reduction. If the local determinant anomaly is \(\alpha Q+L+G\), where \(L_g\,\dd V_g\) is a pointwise conformally invariant density and \(G\) has the stated natural local conformal primitive, integration along the conformal path gives the displayed Polyakov term \emph{plus}
\[
\int_M\omega L_g\,\dd V_g,
\]
the local counterterm difference and the kernel cocycle. A bare divergence decomposition does not by itself provide the needed local primitive. For a nonzero type-B term \(L\), if \(\alpha\ne0\), replacing \(Q\) by the operator-dependent choice \(Q_A=Q+\alpha^{-1}L\) absorbs that linear term while preserving the \(Q\) transformation law. If canonical \(Q\) is fixed, the term must remain unless separately removed by a valid local counterterm. If \(\alpha=0\), it cannot be absorbed by this division. The original question does not specify whether its \(Q\) is fixed or operator dependent.

The aid also reports an exact spectral-power audit. For a nonnegative self-adjoint elliptic \(B\), with zero modes omitted, literal powers satisfy
\[
\zeta_{B^r}(s)=\zeta_B(rs),\qquad
\log\det'(B^r)=r\log\det'B\qquad(r\in\N_{>0}).
\]
No repeated-factor multiplicative anomaly is inserted into this identity. General products or incompatible spectral cuts are different questions. On a flat closed even torus, take the normalized critical operator \(B=P_n=\Delta_+^{n/2}\), with constant kernel, and \(A=B^r\). For constant \(\omega=c\), the reported exact normalization is
\[
-\log\frac{\det' A_{e^{2c}g}}{\det' A_g}=-nrc,
\qquad H(e^{2c}g,g)=-nrc
\]
when the flat local terms use the usual vanishing normalization. Thus the kernel correction cannot be discarded solely because the metric is flat.

The next source problem proposes the natural-density decomposition as sufficient for the determinant program. The present aid gives an audit and conditional integration argument, not an unrestricted theorem proving the exact canonical-\(Q\) formula in every higher even dimension.
\subsection{Short English statement}
Does the low-dimensional Polyakov determinant formula extend to higher even dimensions? The supplied aid identifies the missing type-B term, the choice between canonical and operator-dependent \(Q\), and an exact flat-torus kernel normalization. Literal spectral powers introduce no repeated-factor anomaly.
\subsection{Sources}
\begin{itemize}
\item Problem 29, section extracted from T. Branson and R. Gover, \emph{Origins, applications, and generalizations of the Q-curvature}, in the AIM workshop \emph{Conformal structure in geometry, analysis, and physics}: \url{https://aimath.org/WWN/confstruct/confstruct.pdf}.
\item Local registry \texttt{data/unsolvedmath/problems.json}, permanent ID 20002051, code \texttt{catalogue entry}; canonical \texttt{aim-geometry-notes.json}, \texttt{notes[388]}; updated 14 August 2026. Companion ID 20002050, code \texttt{AIM-GEOMETRY-0388}, \texttt{notes[387]}, appends equation (40) and its operator/determinant assumptions. Following ID 20002052, code \texttt{AIM-GEOMETRY-0390}, \texttt{notes[389]}, contains the referenced natural-density conjecture.
\item The original literature field is empty; no further detailed theorem reference is supplied. Attempt 1: partial result, partially solved status, rigor label proved, novelty confidence low. The machine-generated aid is not peer reviewed and requires independent verification before citation.
\end{itemize}
\end{document}
